\documentclass[12pt]{article}
\usepackage[utf8]{inputenc}
\usepackage{amsmath}
\usepackage{amsfonts}
\usepackage{algorithm}
\usepackage{algorithmic}
\usepackage{booktabs}
\usepackage{graphicx,psfrag,epsf}
\usepackage{enumerate}
\usepackage{natbib}
\usepackage{subcaption}
\usepackage{float}
\renewcommand\bibsection{\section{\refname}}

\usepackage{url}
\usepackage{hyperref}

\addtolength{\oddsidemargin}{-.5in}%
\addtolength{\evensidemargin}{-.5in}%
\addtolength{\textwidth}{1in}%
\addtolength{\textheight}{1.3in}%
\addtolength{\topmargin}{-.8in}%

\def\spacingset#1{\renewcommand{\baselinestretch}%
{#1}\small\normalsize} \spacingset{1}

\newcommand{\thesistitle}{Building mean-variance portfolios that perform in a large universe of assets}
\newcommand{\theauthor}{Kiril Ivanov}
\newcommand{\studentid}{15174433}
\newcommand{\programme}{BSc Econometrics and Data Science}
\newcommand{\institution}{University of Amsterdam}
\newcommand{\department}{Faculty of Economics \& Business}
\newcommand{\supervisor}{Dr. Sanders Barendse}
\newcommand{\thesisdate}{June, 2026}

\hypersetup{
    colorlinks=true,
    linkcolor=black,
    urlcolor=blue,
    pdftitle={\thesistitle},
    pdfauthor={\theauthor},
    citecolor=cyan,
}
\captionsetup[table]{
    width=0.90\textwidth,
}

\captionsetup[figure]{
    width=0.90\textwidth,
    justification=centering
}

\begin{document}
\pagenumbering{gobble}

\begin{titlepage}
    \centering
    \includegraphics[width=0.5\linewidth]{style/logoUvA_en.pdf}\\[2em]
    {\large \institution}\\[0.3em]
    {\normalsize \department}\\[0.3em]
    {\normalsize \programme}\\[2.5em]
    \rule{\linewidth}{0.5pt}\\[0.6em]
    {\LARGE \bfseries \thesistitle}\\[0.4em]
    \rule{\linewidth}{0.5pt}\\[2em]
    {\large \theauthor}\\[0.3em]
    {\normalsize Student ID: \studentid}\\[2.5em]
    {\small
        \begin{tabular}{rl}
            \textit{Supervisor:} & \supervisor \\
        \end{tabular}
    }\\[2.5em]
    {\normalsize \thesisdate}
    \vfill
\end{titlepage}

\newpage
\section*{Statement of Originality}

This document is written by Kiril Raychev Ivanov who declares to take full responsibility for the contents of this document. I declare that the text and the work presented in this document are a reflection of my own efforts and conform to all rules regarding academic integrity. All sources and tools used in its creation, including any AI--based applications, are appropriately acknowledged and referenced. UvA Economics and Business is responsible solely for the supervision of completion of the work and submission, not for the contents.

\newpage
\pagenumbering{arabic}

\begin{abstract}
\noindent Mean-variance optimization often performs poorly out of sample because covariance-matrix inversion can amplify estimation error, producing unstable and concentrated portfolios. This thesis introduces Hierarchical Mean Variance Allocation (HMVA), a long-only portfolio-construction method that targets risk-adjusted performance without directly inverting an estimated covariance matrix. HMVA builds a top-down binary tree by splitting assets to minimize correlation between equal-weighted sub-portfolios, then allocates capital recursively using cluster-level Sharpe proxies based on return estimates. When no positive return estimate is available, it falls back to inverse-volatility bisection. A latent weight filter then blends target weights with prior holdings to reduce turnover when risk estimates are diffuse and to respond more aggressively when return signals shift. HMVA is tested in a fully point-in-time backtest of the top-100 S\&P 500 constituents from July 2000 to January 2024 against mean-variance, minimum-variance, hierarchical risk-parity, and passive benchmarks. HMVA achieves the highest full-sample Sharpe ratio of 0.616, outperforming all benchmarks in full-sample Sharpe ratio on point estimates. Pairwise Sharpe-ratio differences are positive against all comparators in the HMVA hypothesis family, but the studentized bootstrap confidence intervals include zero after multiple-testing correction, so the statistical evidence should be interpreted cautiously.
\end{abstract}


\newpage
\tableofcontents

\newpage
\spacingset{1.45}

\section{Introduction}
\label{sec:introduction}

Mean-variance optimization is one of the central ideas in portfolio theory. In its classical form, it allocates capital by trading off expected return against portfolio variance. The appeal of this framework is clear and intuitive, because if expected returns and covariances were known, an investor could construct an efficient portfolio that offers the best attainable risk-return trade-off. In empirical applications, however, these quantities must be estimated from historical data. This creates a major practical difficulty, because even small errors in expected returns or covariances can translate into large changes in portfolio weights. This problem is especially relevant when the number of assets is large relative to the available estimation window. In that setting, estimated covariance matrices are noisy, expected-return estimates are weak, and classical mean-variance optimization can become unstable out of sample. Instead of using the optimizer to reduce risk, the portfolio may end up amplifying estimation error. As a result, simple allocation rules such as equal weighting often remain difficult benchmarks to beat in realistic out-of-sample tests. This thesis addresses the following question: 

\textit{Can mean and variance estimates be introduced into an allocation framework to outperform both naive and sophisticated benchmarks in risk-adjusted returns out of sample, given identical inputs in a difficult estimation regime?} 

To answer this question, the thesis introduces the Hierarchical Mean Variance Allocation framework designed to retain part of the economic logic of mean-variance optimization while making the mapping from noisy estimates to portfolio weights less direct. Hierarchical Mean Variance Allocation differs from classical mean-variance optimization in three main ways. First, it avoids direct inversion of the estimated covariance matrix. Second, it uses covariance estimates to construct a diversification tree rather than to solve a global quadratic optimization problem. Third, it introduces return information only at the cluster level through recursive risk-adjusted return bisection, so that return estimates affect broader groups of assets before being translated into individual holdings. A latent weight filter then blends the newly estimated target portfolio with previous holdings in order to reduce allocation estimation noise in asset weights.

The empirical analysis evaluates Hierarchical Mean Variance Allocation using a fully point-in-time backtest of the top-100 S\&P 500 constituents from July 2000 to January 2024 with a rolling estimation sample of 126 days. The strategy is compared with passive and active benchmarks, such as equal weight and classical mean-variance optimization portfolios. The analysis also studies performance across market regimes and evaluates the sensitivity of the results to proportional transaction costs. The main empirical finding is that Hierarchical Mean Variance Allocation achieves the highest full-sample point-estimate Sharpe ratio among the strategies considered. The statistical evidence is positive but statistically inconclusive after multiple-testing correction.

\section{Literature review}
\label{sec:lit-review}

\citet{markowitz1952} established that an investor who cares only about the first two moments of portfolio returns should hold a portfolio on the efficient frontier. The classical mean-variance optimal portfolio, which maximizes a scaled trade-off between risk and return subject to the full-investment constraint, lies on this theoretical efficient frontier. In practice, however, the portfolio requires estimates of expected returns and the covariance matrix. Since these inputs are estimated with noise, the resulting allocation can be highly sensitive to small changes in the data. A large literature documents the severity of this estimation problem. \citet{michaud1989} described this failure mode as ``optimization error maximization'' and \citet{chopraziemba1993} further quantified this sensitivity, finding that return errors are approximately ten times more harmful than covariance errors, although both sources of error can cause economically meaningful performance losses. Consistent with this concern, \citet{demiguel2009} found that across fourteen portfolio strategies and seven datasets, sophisticated optimizers did not consistently outperform the naive equal-weighted portfolio once estimation error was taken into account.

One response to this problem is to improve the estimated inputs before optimization, because less noisy inputs are expected to produce more stable portfolios. On the return side, shrinkage ideas date back at least to \citet{stein1956}. These ideas were later adapted for practical portfolio applications in the Black-Litterman framework developed by \citet{blacklitterman1992}. Black-Litterman derived a Bayesian method for combining prior expected returns with investor views under prediction uncertainty. On the covariance side, \citet{ledoitwolf2004} derived an analytically optimal linear shrinkage estimator, and \citet{ledoitwolf2020} later extended this approach by shrinking each eigenvalue individually, improving covariance estimation in high-dimensional settings. Many other approaches to return and covariance regularization and estimation exist across different fields, but they broadly share the same goal. All of these methods aim to make the inputs more stable, so that optimized portfolios become less sensitive to sampling variation.

A related response is to regularize the portfolio weights directly. \citet{frostsavarino1988} were among the first to document that bounded weight constraints can reduce portfolio variance with little loss in expected return. \citet{jagannathanma2003} later formalized this insight by showing that imposing no-short-sale constraints on the minimum-variance portfolio is mathematically equivalent to applying a form of shrinkage to the covariance matrix and \citet{demiguelnogales2009} extended this interpretation by showing that constraints on the weight vector generalize the shrinkage effect and can produce sparser, more stable portfolios that outperform unconstrained sample-based mean-variance optimization.

A more recent response is to redesign the allocation rule so that estimation error is not magnified through covariance-matrix inversion. \citet{lopezdeprado2016} proposed Hierarchical Risk Parity, which builds a similarity hierarchy from the asset correlation matrix and then allocates capital recursively through the tree using variance-based bisection. The appeal of this approach is that the covariance estimate enters the portfolio indirectly as correlations determine a diversification ordering, while allocation proceeds locally rather than through a global inverse. However, it intentionally excludes expected-return information and therefore targets risk allocation rather than mean-variance allocation.

The existing literature therefore motivates a hierarchical method built to use return information without allowing noisy return estimates to dominate the allocation. This thesis develops Hierarchical Mean Variance Allocation as one such approach. The method combines three design principles suggested by the literature. First, it avoids direct covariance inversion. Second, it carefully introduces return estimates only after regularization as well as keeping a long-only profile. Third, it applies a latent weight filter to regularize portfolio weights. In doing so, it attempts to preserve part of the mean-variance objective while reducing the instability usually associated with classical mean-variance optimization.

\section{Methodology}
\label{sec:meth}
\subsection{Estimation} 
The total observation count is denoted by $T$, and the total number of assets is denoted by $N$. Returns are collected in a $T \times N$ return matrix $R$, where each row corresponds to observation index $t$ and each column to asset index $n$, denoted by $R_{t,n}$. The weight assigned by a strategy to asset $n$ at time $t$ is therefore denoted $w_{t,n}$. A time-series window of the form $[t_{\mathrm{start}},t_{\mathrm{end}}]$ is written as $R_{t_{\mathrm{start}}:t_{\mathrm{end}}}$. The rolling training sample size is denoted by $T_h$ and is also referred to as the history length, while the rebalance interval length is denoted by $T_r$. The true covariance matrix and expected return vector are denoted by $\Sigma$ and $\mu$, respectively, and their estimators by $\hat{\Sigma}$ and $\hat{\mu}$. The following estimators are used throughout the thesis for all strategies.

\subsubsection{Exponential Weighting}
Before the return and covariance estimators are applied, a preprocessing is done on the return window using exponential weighting. Following the \cite{riskmetrics1996} approach, let the decay parameter be $\theta = 2^{-1/T_\theta}$, corresponding to a half-life of $T_\theta$ trading days. The original return matrix $R$ is rescaled into the pseudo-return matrix $\tilde{R}$ as:
\begin{equation}
  \tilde{R}_{t} = \sqrt{s_t \, T}\; R_{t}, \qquad
  s_t = \frac{(1-\theta)\,\theta^{\,T-1-t}}{\sum_{i=0}^{T-1}(1-\theta)\,\theta^{\,i}}, 
  \qquad 
  \sum_{t=0}^{T-1} s_t = 1 
  \label{eq:ewma-pseudo}
\end{equation}

The sample moments of $\tilde{R}$ recover the exponentially weighted moments of $R$. This transformation allows the return and covariance estimators to be applied to $\tilde{R}$ as if it were an ordinary return matrix, while automatically inheriting the exponential weighting. Throughout the empirical analysis, the half-life is set equal to the rebalance interval, so $T_\theta=T_r$, matching the decay rate to the forecasting horizon and making the estimates more responsive to the current asset risk-return structure than the corresponding unweighted estimators. By setting the half-life equal to the rebalance interval, information from recent rebalancing periods receives the largest weight, while older observations decay gradually rather than being discarded abruptly. This is useful in a time-varying universe, where correlations, volatilities, and return signals can shift over time, and the allocation should primarily reflect the market structure prevailing near the rebalance date.

\subsubsection{Covariance Estimation}
Random matrix theory, following \citet{marchenkopastur1967}, shows that sample eigenvalues are systematically dispersed away from their population counterparts, with the largest eigenvalues inflated and the smallest eigenvalues deflated. Because portfolio optimizers are especially sensitive to the smallest eigenvalues, this distortion is precisely what destabilizes any allocation that depends on $\hat{\Sigma}$, and correcting the eigenvalues before they enter the allocation is therefore essential.

The analysis adopts the nonlinear shrinkage estimator of \cite{ledoitwolf2012} in the analytical implementation of \cite{ledoitwolf2020}. The idea is to leave the eigenvectors of the sample covariance untouched and to shrink each eigenvalue individually toward the value that minimizes the expected loss against the true covariance. Given the exponential return matrix over the history window $\tilde{R}_{t-T_h:t-1}$, the first step calculates the sample covariance as:
\begin{equation}
\hat{\Sigma}
=
\frac{1}{T_h-1}
\sum_{i=t-T_h}^{t-1}
\left(\tilde{R}_i-\bar{\tilde{R}}\right)^{\top}
\left(\tilde{R}_i-\bar{\tilde{R}}\right),
\qquad
\bar{\tilde{R}}
=
\frac{1}{T_h}
\sum_{i=t-T_h}^{t-1} \tilde{R}_i
\label{eq:sample_cov}
\end{equation}

Shrinkage begins with the eigen-decomposition $\hat{\Sigma} = U\Lambda U^{\top}$, where $\Lambda = \mathrm{diag}(\lambda_1,\dots,\lambda_N)$ is the eigenvalue matrix. Each sample eigenvalue is assigned to its oracle counterpart according to:
\begin{equation}
  \lambda_i^{\star} = \frac{\lambda_i}{\bigl|\,1 - N/T_h - N/T_h \cdot \lambda_i\, \tilde{m}_F(\lambda_i)\,\bigr|^{2}},
  \qquad \tilde{m}_F(\lambda)=\lim_{z \to \lambda + i 0^{+}} \int \frac{1}{t - z}\, dF(t)
  \label{eq:nls}
\end{equation}
where $\tilde{m}_F(\lambda)$ is the boundary value of the Stieltjes transform of the limiting spectral distribution $F$ of $\hat{\Sigma}$, estimated nonparametrically via the kernel density formula of \cite{ledoitwolf2020}. The shrunk covariance estimator is then reconstructed from the original eigenvectors as:
\begin{equation}
  \hat{\Sigma}^{\text{NLS}} = U\,\mathrm{diag}(\lambda_1^{\star},\dots,\lambda_N^{\star})\,U^{\top}
  \label{eq:nls-cov}
\end{equation}
which delivers the asymptotically optimal correction under Frobenius loss within the class of rotation-equivariant estimators, that is, among all estimators that share the sample eigenvectors and adjust only the eigenvalues. From now on it is assumed that referring to a covariance estimator $\hat{\Sigma}$ is equivalent as referring to $\hat{\Sigma}^{\text{NLS}}$.

\subsubsection{Return Estimation}
To estimate the return signal, the analysis uses the return matrix over the history window $\tilde{R}_{t-T_h:t-1}$ and a covariance estimate $\hat{\Sigma}$ computed over the same return window. The first step constructs the views. To avoid the reversal effects in short horizons, as proposed in \cite{jegadeesh1993}, the analysis skips the most recent $T_s$ days and sets $T_s=T_r$ to exclude returns only since the previous rebalance of the portfolio. Thus averaging returns over the interval $[t-T_h,t-T_r]$ is equivalent to the first moment of $\tilde{R}_{t-T_h:t-T_r}$:
\begin{equation}
    Q = (T_h - T_r)^{-1}\sum_{t-T_h < i < t-T_r} \tilde{R}_i
    \label{eq:view_const}
\end{equation}

Since the model forms an asset-specific view for each asset, the view vector $Q$ is a dense $1 \times N$ vector and the view matrix is $P=I_N$. Thus each view is an absolute view on one asset, rather than a relative view between assets. The view uncertainty is set equal to the diagonal covariance of each asset, so $\Omega=\mathrm{diag}(\hat{\Sigma})$. The prior return is set to the cross-sectional average return over the estimation window, rounded with an upward bias to the second decimal percentage point to encourage investment, and assigned equally to every asset:
\begin{equation}
    \pi
    =
    \mathbf{1}_N
    \cdot
    \left(
    \left\lceil
    100 \cdot (T_h N)^{-1}
    \sum_{n=1}^{N}\sum_{i=t-T_h}^{t-1} \tilde{R}_{i,n}
    \right\rceil / 100
    \right)
    \label{eq:prior_bl}
\end{equation}
The posterior return estimate then follows from \cite{blacklitterman1992}:
\begin{equation}
    \hat{\mu}
    =
    \bigl((\tau\hat{\Sigma})^{-1} + P^{\top}\Omega^{-1}P\bigr)^{-1}
    \bigl((\tau\hat{\Sigma})^{-1}\pi + P^{\top}\Omega^{-1}Q\bigr)
    \label{eq:BL_formula}
\end{equation}
When $P=I_N$, $\Omega=\mathrm{diag}(\hat{\Sigma})$, and $\tau=1$, this can be written as:
\begin{equation}
    \hat{\mu}
    =
    \pi
    +
    \hat{\Sigma}
    \bigl(\hat{\Sigma}+\mathrm{diag}(\hat{\Sigma})\bigr)^{-1}
    (Q-\pi)
    \label{eq:BL_simplified}
\end{equation}

The final return signal is therefore a Bayesian point estimate that blends a conservative cross-sectional prior with the first moment of sample returns, producing a natural shrinkage toward a nonlinear target.

\subsection{Hierarchical Mean Variance Allocation}
\label{sec:hmva}

Hierarchical Mean Variance Allocation (HMVA) constructs an investment decision in three allocation stages after the return vector and covariance matrix have been estimated. First, it constructs a top-down binary asset hierarchy by recursively separating clusters whose equal-weighted returns have low cross-correlation. Second, it allocates capital through the hierarchy using cluster-level Sharpe approximations. Third, it applies a latent weight filter that smooths the target allocation between rebalances.

After the return vector $\hat{\mu}$ and covariance matrix $\hat{\Sigma}$ have been estimated from the training sample, these estimates are used to build the hierarchical allocation. Starting with the full asset universe as the root cluster, HMVA recursively splits each parent cluster into two disjoint child clusters, denoted $L$ and $R$. For each candidate split, define the equal-weighted return of a cluster $C \in \{L, R \}$ as the return on a portfolio that assigns equal weight $1/n_C$ to each asset in the cluster, where $n_C=|C|$. The estimated variance and expected return of this equal-weighted cluster portfolio are:
\begin{equation}
\widehat{\mathrm{Var}}_{\mathrm{EW}}(C)
=
n_C^{-2}
\sum_{i \in C,\, j \in C}
\hat{\Sigma}_{ij}, 
\qquad
\hat{\mu}_{\mathrm{EW}}(C)
=
n_C^{-1}
\sum_{i \in C}
\hat{\mu}_{i}
\label{eq:obj_mu_sd}
\end{equation}
and estimated covariance between the equal-weighted child portfolios is:
\begin{equation}
\widehat{\mathrm{Cov}}_{\mathrm{EW}}(L,R)
=
\frac{
\sum_{i\in L,\,j\in R}\hat{\Sigma}_{ij}
}{
n_L n_R
}
\label{eq:cluster_covariance}
\end{equation}
HMVA chooses the split that minimizes the estimated correlation between the two equal-weighted child portfolios:
\begin{equation}
f(L,R)
=
\widehat{\mathrm{Corr}}_{\mathrm{EW}}(L,R)
=
\frac{
\widehat{\mathrm{Cov}}_{\mathrm{EW}}(L,R)
}{
\sqrt{
\widehat{\mathrm{Var}}_{\mathrm{EW}}(L)
\widehat{\mathrm{Var}}_{\mathrm{EW}}(R)
}
}
=
\frac{
\sum_{i\in L,\,j\in R}\hat{\Sigma}_{ij}
}{
n_L n_R
\sqrt{
\widehat{\mathrm{Var}}_{\mathrm{EW}}(L)
\widehat{\mathrm{Var}}_{\mathrm{EW}}(R)
}
}
\label{eq:split-obj}
\end{equation}

The motivation for this objective follows from the variance decomposition of a two-branch portfolio. Suppose a fraction $\alpha$ is allocated to branch $L$ and $1-\alpha$ to branch $R$. Let
$
\hat{\sigma}_L^2=\widehat{\mathrm{Var}}_{\mathrm{EW}}(L), \text{ }
\hat{\sigma}_R^2=\widehat{\mathrm{Var}}_{\mathrm{EW}}(R), \text{ }
\hat{\rho}_{LR}=\widehat{\mathrm{Corr}}_{\mathrm{EW}}(L,R).
$
The variance of the parent portfolio is then:
\begin{equation}
\widehat{\mathrm{Var}}(Parent)
=
\alpha^2 \hat{\sigma}_L^2
+
(1-\alpha)^2 \hat{\sigma}_R^2
+
2\alpha(1-\alpha)\hat{\rho}_{LR}\hat{\sigma}_L\hat{\sigma}_R
\label{eq:branch_variance_decomposition}
\end{equation}

For fixed branch volatilities and a fixed subsequent allocation rule, the cross-branch correlation is the component of the split that directly controls the diversification benefit. Since:
\begin{equation}
\frac{
\partial \widehat{\mathrm{Var}}(Parent)
}{
\partial \hat{\rho}_{LR}
}
=
2\alpha(1-\alpha)\hat{\sigma}_L\hat{\sigma}_R
\geq 0
\label{eq:variance_corr_derivative}
\end{equation}
for any interior allocation $\alpha\in(0,1)$, reducing cross-branch correlation weakly reduces parent-level variance. The same intuition becomes especially transparent under inverse-volatility bisection. If capital is allocated using inverse-volatility weights
\[
\alpha_L
=
\frac{1/\hat{\sigma}_L}{1/\hat{\sigma}_L+1/\hat{\sigma}_R}
=
\frac{\hat{\sigma}_R}{\hat{\sigma}_L+\hat{\sigma}_R},
\qquad
\alpha_R
=
\frac{\hat{\sigma}_L}{\hat{\sigma}_L+\hat{\sigma}_R}
\]
then substituting these weights into \eqref{eq:branch_variance_decomposition} gives:
\begin{equation}
\widehat{\mathrm{Var}}(Parent)
=
\frac{
2\hat{\sigma}_L^2\hat{\sigma}_R^2(1+\hat{\rho}_{LR})
}{
(\hat{\sigma}_L+\hat{\sigma}_R)^2
}
\label{eq:ivp_variance_correlation}
\end{equation}

Conditional on the two branch volatilities, parent variance is therefore increasing in $\hat{\rho}_{LR}$. Minimizing the equal-weighted branch correlation is consequently aligned with constructing a hierarchy whose upper-level branches are maximally diversifying. Volatility and expected return information are then introduced in the recursive allocation step. This creates a separation of roles where the tree determines the diversification structure, while the allocation rule determines how much capital each branch receives to exploit the return signal.

Once the tree is constructed, capital is allocated top-down. At each internal node with left and right children $L$ and $R$, equal-weighted cluster Sharpe approximations are formed to reduce signal-estimation noise through cluster averaging as:
\begin{equation}
  \widehat{SR}(C) 
  =
  \max\!\left(
    \frac{\hat{\mu}_\mathrm{EW}(C)}
         {\sqrt{\widehat{\mathrm{Var}}_{\mathrm{EW}}(C)}},
    0
  \right)
  =
  \max\!\left(
    \frac{
    n_C^{-1}\sum_{i \in C}\hat{\mu}_{i}
    }{
    \sqrt{n_C^{-2}\sum_{i,j \in C}\hat{\Sigma}_{ij}}
    },
    0
  \right),
  \qquad C\in\{L,R\}
  \label{eq:sharpe-proxy}
\end{equation}
The numerator is the equal-weighted expected return implied by $\hat{\mu}$, while the denominator is the equal-weighted volatility implied by $\hat{\Sigma}$. The truncation at zero prevents negative return estimates from receiving negative or short allocations, preserving the long-only nature of the strategy to ensure stability according to \cite{jagannathanma2003}. If both child clusters have zero Sharpe proxies, the allocation reverts to inverse-volatility bisection:
\begin{equation}
  \widehat{IV}(C)
  =
  \frac{1}{\sqrt{\widehat{\mathrm{Var}}_{\mathrm{EW}}(C)}}
  =
  \frac{
  1
  }{
  \sqrt{n_C^{-2}\sum_{i,j \in C}\hat{\Sigma}_{ij}}
  }
  \label{eq:inv_var}
\end{equation}
When no return estimate is available, every split is forced to use this inverse-volatility branch. This risk-only variant is referred to as the minimum-variance version of HMVA, denoted HMVA-mv. The allocation to the left child is therefore:
\begin{equation}
\alpha_L 
=
\frac{\widehat{SR}(L)}{\widehat{SR}(L)+\widehat{SR}(R)},
\qquad 
\alpha_R = 1-\alpha_L
\quad \text{if } \widehat{SR}(L)+\widehat{SR}(R)>0
\end{equation}
and otherwise:
\begin{equation}
\alpha_L 
=
\frac{\widehat{IV}(L)}{\widehat{IV}(L)+\widehat{IV}(R)},
\qquad 
\alpha_R = 1-\alpha_L
\label{eq:sharpe-alloc}
\end{equation}
Thus, when return information is positive, the branch with the higher cluster-level risk-adjusted expected return receives more capital. When return information is absent or non-positive, the rule falls back to a risk-based allocation. Because capital is allocated recursively through a binary tree, the hierarchy also imposes implicit concentration control without requiring explicit upper bounds on individual asset weights.

To curb portfolio changes generated by estimation error in $\hat{\mu}$ and $\hat{\Sigma}$, while still retaining relevant new information, HMVA interprets rebalancing as a linear filtering problem for the latent portfolio weights. Let $w_t^{\ast}$ denote the unobserved target allocation that would be optimal under the true investment opportunity set at rebalance date $t$. The allocation produced by the HMVA tree before turnover smoothing, denoted by $w_{\mathrm{new}}$, is treated as a noisy measurement of this latent target. The previously held allocation, denoted by $w_{\mathrm{old}}$, is treated as the prior estimate of the same latent target carried forward from the preceding rebalance. This interpretation follows the general state-space view of the Kalman filter, in which an unobserved state is updated recursively by combining a prior estimate with a noisy observation, as in \cite{kalman1960}. This gives the following local state-space representation:
\begin{equation}
w_t^{\ast} = w_{t-1}^{\ast} + \eta_t, 
\qquad
w_{\mathrm{new}} = w_t^{\ast} + \nu_t
\label{eq:state_innov} 
\end{equation}
where $\eta_t$ is the state innovation and $\nu_t$ is the measurement error in the newly estimated target portfolio. The first equation states that the latent optimal allocation evolves gradually over time. The second equation states that the newly constructed HMVA target is informative, but contaminated by the estimation error in $\hat{\mu}$ and $\hat{\Sigma}$. Under this interpretation, turnover control is equivalent to filtering a noisy portfolio-weight signal rather than mechanically trading all the way to the latest target.

Under the standard Kalman filter update, the posterior estimate is a weighted average of the prior state estimate and the new measurement, with the weights determined by their relative uncertainty, as shown by \cite{kalman1960}. In the scalar-gain version used here, the smoothed allocation is therefore written as:
\begin{equation}
\tilde{w} = K w_{\mathrm{new}} + (1-K) w_{\mathrm{old}} \label{eq:kalman_update}
\end{equation}
where $K \in [0,1]$ is the Kalman gain. A high value of $K$ places more weight on the new target, while a low gain shrinks the allocation toward the previous holdings. Thus, the Kalman update has the same practical form as a turnover filter as it implements partial adjustment toward the frictionless target and reduces unnecessary trading when the new target is likely to be dominated by estimation noise. This is also consistent with dynamic portfolio-choice models under proportional transaction costs, where full rebalancing is generally replaced by no-trade regions or partial-adjustment policies \cite{constantinides1986} and \cite{demiguelmei2016}, and with the view that turnover penalization can act as a form of regularization under model uncertainty \cite{hautschvoigt2019}.

To calibrate this gain, HMVA defines a signal-based measure of process uncertainty:
\begin{equation}
Y =   \frac{ \sqrt{ \sum_{\text{all n}} \left( \hat{\mu}_{\text{new},n} - \hat{\mu}_{\text{old},n} \right)^2 } }{ \sqrt{ \sum_{\text{all n}}  \hat{\mu}_{\text{old},n} ^2 } + \epsilon }\label{eq:Y}
\end{equation}
where $\epsilon > 0$ is included for numerical stability. A large value of $Y$ indicates that the estimated return environment has changed materially, so the latent optimal allocation is more likely to have moved away from the previous holdings. In Kalman-filter terminology, $Y$ therefore plays the role of a proxy for process uncertainty. This choice is motivated by the fact that expected-return estimates are a major source of portfolio instability, so a large change in $\hat{\mu}$ should increase the willingness to update the allocation \cite{chopraziemba1993}. When the filtering is used in the risk-only variant, then $Y=1$ always. Similarly, a measure of measurement uncertainty is derived from the eigenvalue structure of $\hat{\Sigma}$. The entropy-based uncertainty proxy is then defined as:
\begin{equation}
X = -\frac{ \sum_{i=1}^{N} p_i \log p_i }{ \log N }, \quad p_i = \frac{ \lambda_i }{ \sum_{\text{all n}} \lambda_n } \label{eq:X}
\end{equation}
where $\lambda_i$ for $i \in \{1,\dots,N\}$ are the eigenvalues of the covariance matrix. The denominator $\log N$ rescales the entropy onto $[0,1]$. The lower bound is approached when a single eigenvalue dominates the spectrum, indicating a sharply concentrated risk structure. The upper bound is attained when all eigenvalues are equal, indicating a maximally diffuse covariance structure. This normalized spectral entropy follows the entropy-based diversification diagnostics of \cite{meucci2009}. A high value of $X$ therefore indicates that the estimated risk structure is less sharply identified, so the newly measured target allocation should be trusted less. In the filtering interpretation, $X$ acts as a proxy for measurement uncertainty. The Kalman gain is then set as a normalized function of these two uncertainty proxies:
\begin{equation}
K = \frac{Y}{X+Y} \label{eq:gain}
\end{equation}

When the latent investment opportunity set appears to have changed substantially, $Y$ is large relative to $X$, and the gain increases. The portfolio then moves more aggressively toward the newly estimated HMVA target. Conversely, when the covariance estimate is diffuse and the return signal is relatively stable, $X$ is large relative to $Y$, and the gain decreases. The portfolio then remains closer to its previous holdings. The resulting change in portfolio weights induced by the filter is therefore:
\begin{equation}
\Delta w = K \left( w_{\mathrm{new}} - w_{\mathrm{old}} \right) \label{eq:turnover}
\end{equation}

Lower values of $K$ therefore imply lower turnover, while higher values of $K$ imply faster adjustment. Under the linear state-space model in \eqref{eq:state_innov}, the Kalman update has an optimal linear-estimation interpretation when the relevant process and measurement uncertainty are correctly specified \cite{kalman1960}. In the present application, $X$ and $Y$ are not estimated noise variances, but proportional scalar proxies for measurement and process uncertainty. Therefore, the update in \eqref{eq:kalman_update} is best interpreted as a scalar linear estimator of the latent target weights, rather than as the exact optimal Kalman filter. After the update, the smoothed weights $\tilde{w}$ are renormalized to sum to one. A complete pseudo-code implementation of HMVA and HMVA-mv, along with a weblink to the codebase, is provided in Appendix~\ref{app:pseudo-code}$\,$.

\subsection{Experiment}
\label{sec:data}
\subsubsection{Data and Design}
The empirical experiment uses daily CRSP data with four primary fields: \texttt{PERMNO} (unique stock identifier), \texttt{DlyCalDt} (calendar date), \texttt{DlyClose} (split-adjusted closing price), \texttt{DlyRet} (daily total return), and daily market capitalization. To populate the return matrix, \texttt{DlyRet} is used directly when available for a given stock-date pair. When it is missing, returns are computed from consecutive split-adjusted closing prices. Finally, public NBER recession data are used to construct a dummy variable indicating whether day $t$ falls in a recession. The rebalance window is monthly, hence $T_r=21$, and the estimation history is six months, thus $T_h=126$ in trading days. This rolling sample size is chosen to provide a difficult estimation regime of $N/T = 100/126 \approx 0.8\,$, given the scope of the research question and the rebalance schedule is standard for portfolio performance evaluation.

\subsubsection{Point-in-time Universe}
\label{sec:data-pit}
S\&P 500 membership is reconstructed from a historical constituents file that records membership in range format, such that each row identifies a PERMNO and the interval during which it was an S\&P 500 constituent. This mitigates survivorship bias, since only stocks that were members of the index on the specific rebalance date are eligible for inclusion in the portfolio. At each $T_r$-day rebalance date $t$, approximately monthly, the investable universe is constructed by first finding the set of PERMNOs with active S\&P 500 membership on $t$ from the constituents file, then restricting it to the top-100 by market capitalization as of $t-1$, using the most recent available daily cap figure. The analysis then retains only those PERMNOs with a complete return data profile over the full estimation history window $[t - T_h, t - 1]$. Stocks that fail this criterion are not eligible for allocation at that rebalance. Operationally, their portfolio weights are set to zero and the remaining eligible stocks are renormalized to sum to one.

The resulting universe changes at every rebalance date, reflecting actual index additions, deletions, and corporate events in real time. The backtest period spans July 3, 2000 to January 1, 2024, with the first usable rebalance date occurring after the six-month warm-up period ($T_h=126$) from the data start of January 1, 2000.

\subsubsection{Benchmark Strategies}
\label{sec:data-strategies}
The empirical comparison includes passive benchmarks and active allocation rules. The two passive benchmarks are the equal-weighted portfolio (EW) and the market-capitalization-weighted portfolio over the top-100 universe (SPY-100). The active strategies are divided into risk-only methods and return-based methods. The first group includes HMVA-mv, Hierarchical Risk Parity (HRP), and Global Minimum Variance (GMV). The second includes Mean Variance Optimization (MVO) and HMVA. Further details on the benchmark strategies are provided in Appendix~\ref{app:benchmarks}. Every strategy that requires estimation draws on an identical information set of the history window of length $T_h=126$ and uses the monthly rebalance interval $T_r=21$ as defined above, applied to the same point-in-time universe of the top-100 stocks. The holding window, rebalance schedule and universe are fixed across strategies to ensure that any difference in realized performance is attributable to the construction mechanism rather than to the data each strategy is allowed to see. 

\subsubsection{Turnover and Transaction-cost}
\label{sec:turnover-costs}

The baseline empirical comparison is conducted without transaction costs in order to isolate the allocation performance of the competing strategies. However, because the strategies differ substantially in how actively they rebalance, turnover and transaction costs are also recorded and evaluated as a robustness exercise. 

At each rebalance date, turnover is measured as the change required to move from the drifted pre-trade portfolio to the newly chosen target portfolio. Let $w_{t-T_r,n}$ denote the portfolio weight assigned to asset $n$ at the previous rebalance date, and let $R_{i,n}$ denote the realized return of asset $n$ between rebalances. Immediately before the next rebalance at time $t$, the previous portfolio has drifted mechanically with realized returns. The drifted pre-trade weight of asset $n$ is therefore:
\begin{equation}
w^{\mathrm{drift}}_{t,n}
=
\frac{
w_{t-T_r,n}
\prod_{i=t-T_r}^{t-1}(1+R_{i,n})
}{
\sum_{m}
w_{t-T_r,m}
\prod_{i=t-T_r}^{t-1}(1+R_{i,m})
}.
\label{eq:drifted_weight_design}
\end{equation}

Let $w_{t,n}$ denote the post-rebalance target weight assigned to asset $n$ at rebalance date $t$. The turnover at this rebalance is then:
\begin{equation}
\mathrm{TO}_t
=
\sum_{n}
\left|
w_{t,n}
-
w^{\mathrm{drift}}_{t,n}
\right|.
\label{eq:turnover_design}
\end{equation}

Because the investable universe is time-varying, the weight vectors at consecutive rebalances may not contain the same assets. Before computing turnover, both $w_t$ and $w_t^{\mathrm{drift}}$ are therefore aligned on the union of assets appearing in either vector. Assets that are absent from one of the two vectors are assigned weight zero. This ensures that newly opened positions and fully liquidated positions are both counted in turnover. Transaction costs are modeled as proportional linear costs charged on turnover. Let $\kappa$ denote the transaction-cost rate per unit of traded capital. The cost deducted at rebalance date $t$ is:
\begin{equation}
c_t
=
\kappa \, \mathrm{TO}_t
=
\kappa
\sum_n
\left|
w_{t,n}
-
w^{\mathrm{drift}}_{t,n}
\right|.
\label{eq:transaction_cost_design}
\end{equation}

The after-cost portfolio return over the holding interval following rebalance date $t$ is computed by subtracting this cost from the corresponding pre-cost portfolio return:
\begin{equation}
r^{\mathrm{net}}_{t}
=
r^{\mathrm{gross}}_{t}
-
c_t .
\label{eq:net_return_design}
\end{equation}
and mean turnover is then reported as the average turnover per rebalance:
\begin{equation}
\overline{\mathrm{TO}}
=
\frac{1}{\# \text{rebalances}}
\sum_{t \in \text{rebalances}}
\mathrm{TO}_t.
\label{eq:mean_turnover_design}
\end{equation}

This statistic measures how much capital is traded on average at each monthly rebalance. The cost model is intentionally simple, similar to the cost model in \cite{lobofazelboyd2007}, and is only used as a robustness check. It does not include nonlinear market impact, time-varying bid-ask spreads, liquidity constraints, or capacity limits.

\subsubsection{Statistical Inference}
\label{inf_stats}
The research question motivates two pre-specified hypothesis families, corresponding to the return-based and risk-only versions of the proposed allocation framework. The first family tests HMVA against the return-based active comparator MVO and the passive comparators EW and SPY-100. The second family tests HMVA-mv against the risk-only active comparators GMV and HRP, as well as the passive comparators EW and SPY-100. These hypotheses evaluate whether the proposed hierarchical allocation improves out-of-sample risk-adjusted performance relative to both classical sophisticated and naive allocation rules given identical estimation inputs. 

Accordingly, since a risk-free rate is not included, the reported Sharpe ratios are return-to-volatility ratios based on raw portfolio returns rather than excess-return Sharpe ratios. Following the bootstrap-based approach of \cite{ledoitwolf2008}, pairwise differences in Sharpe ratios are evaluated using a studentized time-series bootstrap. For two strategies with realized return series $\{r_{t,1}\}_{t=1}^{T}$ and $\{r_{t,2}\}_{t=1}^{T}$, the null hypothesis is:
\begin{equation}
    H_0:\mathrm{SR}_1-\mathrm{SR}_2=0
\end{equation}
The observed annualized Sharpe-ratio difference is:
\begin{equation}
\widehat{\Delta \mathrm{SR}}
=
\widehat{\mathrm{SR}}_1-\widehat{\mathrm{SR}}_2,
\qquad
\widehat{\mathrm{SR}}_i
=
\frac{\bar r_i}{\hat \sigma_i}\sqrt{252}
\label{eq:LW_Sharpe_diff}
\end{equation}
where $\bar r_i$ and $\hat\sigma_i$ are the sample mean and standard deviation of strategy $i$.

Serial dependence and non-normality are accounted for by resampling paired observations $(r_{t,1},r_{t,2})$ with a circular block bootstrap. Blocks have length $T_r$, equal to one rebalance interval, and are wrapped circularly at the end of the sample. This preserves both within-strategy time-series dependence and cross-strategy dependence. For each bootstrap replication $b=1,\dots,B$, a pseudo-sample of length $T$ is drawn and the bootstrap Sharpe-ratio difference is computed as:
\begin{equation}
\widehat{\Delta \mathrm{SR}}_b^{*}
=
\widehat{\mathrm{SR}}_{1,b}^{*}
-
\widehat{\mathrm{SR}}_{2,b}^{*}
\label{eq:LW_boot_diff}
\end{equation}

The statistic is studentized using a heteroskedasticity-and-autocorrelation consistent delta-method standard error. The long-run covariance terms entering this standard error are estimated using the \cite{neweywest1987} estimator with lag length $T_r-1$ chosen according to \cite{neweywest1994}. The standard error of $\widehat{\Delta \mathrm{SR}}$ is denoted by $\widehat{s}$, and the corresponding bootstrap standard error by $\widehat{s}_b^{*}$. The studentized bootstrap statistic is then:
\begin{equation}
t_b^{*}
=
\frac{
\widehat{\Delta \mathrm{SR}}_b^{*}
-
\widehat{\Delta \mathrm{SR}}
}{
\widehat{s}_b^{*}
},
\qquad b=1,\dots,B
\end{equation}
Studentized bootstrap confidence intervals are also reported. If $q_{\alpha/2}$ and $q_{1-\alpha/2}$ denote the empirical quantiles of $\{t_b^{*}\}_{b=1}^{B}$, the $(1-\alpha)$ confidence interval is:
\begin{equation}
\text{C.I. =}
\left[
\widehat{\Delta \mathrm{SR}}
-
q_{1-\alpha/2}\widehat{s},
\quad
\widehat{\Delta \mathrm{SR}}
-
q_{\alpha/2}\widehat{s}
\right]
\label{eq:LW_studentized_CI}
\end{equation}
and the two-sided bootstrap p-value is computed as:
\begin{equation}
\hat p
=
B^{-1}
\sum_{b=1}^{B}
\mathbf{1}
\left[
|t_b^{*}|
\geq
|\frac{\widehat{\Delta \mathrm{SR}}}{\widehat{s}}|
\right]
\label{eq:LW_studentized_p}
\end{equation}

Equality of Sharpe ratios is rejected at level $\alpha$ when zero lies outside this interval. In the empirical analysis, $B=2{,}000$ bootstrap replications are used, the block length is set to $T_r=21$, and 95\% confidence intervals are reported, where the replications and block length are suggested by \cite{ledoitwolf2008}. Because several pairwise Sharpe-ratio comparisons are evaluated, raw $p$-values can overstate the strength of the evidence against the null when interpreted within a family of related tests. The empirical analysis therefore applies a multiple-testing correction within the pre-specified hypothesis families. 

Within each family, the Holm-Bonferroni procedure, derived in \cite{holm1979}, is used to control the family-wise error rate. Let $p_1,\dots,p_m$ denote the raw $p$-values in a family of $m$ related tests, sorted in increasing order as:
\begin{equation}
p_{(1)}\le p_{(2)}\le \cdots \le p_{(m)}
\end{equation}
The Holm procedure rejects the first ordered hypothesis if:
\begin{equation}
p_{(1)} \le \frac{\alpha}{m}
\end{equation}
and then continues sequentially, rejecting the $k$-th ordered hypothesis only if:
\begin{equation}
p_{(k)} \le \frac{\alpha}{m-k+1}
\label{eq:holm-critical}
\end{equation}

Once a hypothesis is not rejected, all remaining hypotheses in that family are also not rejected.
The corresponding adjusted $p$-values are computed as:
\begin{equation}
p^{\mathrm{Holm}}_{(k)}
=
\max_{j\le k}
\left\{
(m-j+1)p_{(j)}
\right\}
\label{eq:holm-adj}
\end{equation}
with the resulting sequence truncated above at one. The adjustment is applied within each hypothesis family rather than across all reported comparisons jointly, because the families answer distinct economic questions and are interpreted separately. Consequently, rejection within one family is not interpreted as evidence for a global claim of superiority across all benchmarks.


\section{Empirical Results}
\label{sec:empirical-results}

\subsection{Backtest Results}
\label{sec:res-full}

Table~\ref{tab:res-main} reports the full-sample performance of the selected strategy representatives, the two proposed hierarchical strategies, and the two passive benchmarks over the evaluation window with zero transaction costs. The table reports annualized return, annualized volatility, Sharpe ratio, Sortino ratio, maximum drawdown, and the mean turnover. HMVA records the highest Sharpe ratio at $0.616$ in the full-sample comparison and HMVA-mv records the second-highest Sharpe ratio, at $0.581$, followed by GMV, HRP, EW, SPY-100, and MVO.

\begin{table}[H]
\centering
\caption{Full-sample performance of the best representative of each family alongside the hierarchical strategies and passive benchmarks, CRSP point-in-time top-100 S\&P 500, July 2000 to January 2024, $T_h=126$, monthly rebalancing, zero transaction costs. MDD is the maximum drawdown, and $\overline{\mathrm{TO}}$ is the mean monthly turnover.}
\label{tab:res-main}
\begin{tabular}{lrrrrrr}
\toprule
Strategy & Ann.\ Ret. & Ann.\ Vol. & Sharpe & Sortino & MDD & $\overline{\mathrm{TO}}$ \\
\midrule
HMVA    & 10.9\% & 17.7\% & 0.616 & 0.801 & $-$31.1\% & 0.782 \\
MVO     &  4.7\% & 31.1\% & 0.151 & 0.185 & $-$69.6\% & 1.490 \\
\midrule
HMVA-mv &  8.7\% & 14.9\% & 0.581 & 0.745 & $-$30.9\% & 0.670 \\
GMV     &  7.4\% & 14.0\% & 0.532 & 0.680 & $-$38.1\% & 0.886 \\
HRP     &  7.9\% & 15.8\% & 0.498 & 0.627 & $-$43.5\% & 0.400 \\
\midrule
EW      &  6.6\% & 19.3\% & 0.345 & 0.437 & $-$55.4\% & 0.101 \\
SPY-100 &  6.3\% & 19.3\% & 0.328 & 0.419 & $-$55.2\% & 0.027 \\
\bottomrule
\end{tabular}
\end{table}

\begin{figure}[!ht]
\centering
\includegraphics[width=0.9\linewidth]{results/backtest/base/equity_curves_log.png}
\caption{Cumulative log wealth index for all strategies over the full period}
\end{figure}

The return-based active benchmark has a lower full-sample Sharpe ratio than HMVA. MVO records a Sharpe ratio of $0.151$ and also has the highest annualized volatility in the table, at $31.1\%$, and the largest maximum drawdown, at $69.6\%$. Among the risk-only strategies, HMVA-mv records a Sharpe ratio of $0.581$, compared with $0.532$ for GMV and $0.498$ for HRP. HMVA-mv has an annualized volatility of $14.9\%$, while GMV and HRP have annualized volatilities of $14.0\%$ and $15.8\%$, respectively. The maximum drawdown of HMVA-mv is $30.9\%$, compared with $38.1\%$ for GMV and $43.5\%$ for HRP. The two passive benchmarks have similar full-sample performance. EW records a Sharpe ratio of $0.345$, while SPY-100 records a Sharpe ratio of $0.328$. Their maximum drawdowns are $55.4\%$ and $55.2\%$, respectively. Compared with these benchmarks, HMVA and HMVA-mv both record a higher Sharpe ratio and a smaller maximum drawdown. 


\subsection{Statistical Results}
\label{sec:res-tests}

Tables~\ref{tab:res-tests-hmva} and~\ref{tab:res-tests-hmva-mv} report pairwise Sharpe-ratio tests for the two pre-specified hypothesis families. The first family compares HMVA with MVO, EW, and SPY-100. The second family compares HMVA-mv with GMV, HRP, EW, and SPY-100.

\begin{table}[!ht]
\centering
\caption{Pairwise tests of HMVA against the comparators in the HMVA hypothesis family. The Sharpe difference is computed as HMVA minus the comparator. LW C.I. is the $95\%$ studentized Ledoit-Wolf bootstrap confidence interval, LW $p$ is the raw two-sided bootstrap $p$-value, and Holm $p$ is the multiple-testing-corrected $p$-value within this hypothesis family.}
\label{tab:res-tests-hmva}
\begin{tabular}{lrrrr}
\toprule
Comparator & SR diff & LW C.I. & LW $p$ & Holm $p$\\
\midrule
MVO  & $+0.369$ & $[-0.005, +0.766]$ & 0.059 & 0.154\\
EW      & $+0.243$ & $[-0.087, +0.575]$ & 0.154 & 0.154\\
SPY-100 & $+0.259$ & $[-0.072, +0.586]$ & 0.133 & 0.154\\
\bottomrule
\end{tabular}
\end{table}

In the HMVA hypothesis family, all reported Sharpe-ratio differences are positive. The largest difference is against MVO, where HMVA exceeds the comparator by $+0.369$. The corresponding 95\% confidence interval is $[-0.005,+0.766]$, with a raw bootstrap $p$-value of $0.059$ and a Holm-adjusted $p$-value of $0.154$. Against EW, the Sharpe-ratio difference is $+0.243$, with a confidence interval of $[-0.087,+0.575]$, a raw $p$-value of $0.154$, and a Holm-adjusted $p$-value of $0.154$. Against SPY-100, the Sharpe-ratio difference is $+0.259$, with a confidence interval of $[-0.072,+0.586]$, a raw $p$-value of $0.133$, and a Holm-adjusted $p$-value of $0.154$. In all three comparisons, the reported confidence interval contains zero.

\begin{table}[!ht]
\centering
\caption{Pairwise tests of HMVA-mv against the comparators in the HMVA-mv hypothesis family. The Sharpe difference is computed as HMVA-mv minus the comparator. LW C.I. is the $95\%$ studentized Ledoit-Wolf bootstrap confidence interval, LW $p$ is the raw two-sided bootstrap $p$-value, and Holm $p$ is the multiple-testing-corrected $p$-value within this hypothesis family.}
\label{tab:res-tests-hmva-mv}
\begin{tabular}{lrrrr}
\toprule
Comparator & SR diff & LW C.I. & LW $p$ & Holm $p$\\
\midrule
GMV  & $+0.049$ & $[-0.137, +0.231]$ & 0.603 & 0.603\\
HRP  & $+0.073$ & $[-0.157, +0.310]$ & 0.541 & 0.603\\
EW   & $+0.202$ & $[-0.090, +0.507]$ & 0.188 & 0.376\\
SPY-100 & $+0.217$ & $[-0.089, +0.529]$ & 0.173 & 0.376\\
\bottomrule
\end{tabular}
\end{table}

In the HMVA-mv hypothesis family, all reported Sharpe-ratio differences are also positive. Relative to GMV, HMVA-mv records a Sharpe-ratio difference of $+0.049$, with a confidence interval of $[-0.137,+0.231]$, a raw $p$-value of $0.603$, and a Holm-adjusted $p$-value of $0.603$. Relative to HRP, the difference is $+0.073$, with a confidence interval of $[-0.157,+0.310]$, a raw $p$-value of $0.541$, and a Holm-adjusted $p$-value of $0.603$. Against the passive benchmarks, HMVA-mv records a difference of $+0.202$ relative to EW and $+0.217$ relative to SPY-100. The corresponding Holm-adjusted $p$-values are $0.376$ in both cases. Each confidence interval in this family contains zero as well.


\subsection{Performance across regimes}
\label{sec:res-regime}

This subsection reports performance separately for crisis and calm regimes. Crisis periods are defined using the NBER recession windows that fall within the evaluation sample. All remaining observations are classified as calm periods. Table~\ref{tab:crisis-periods-combined} lists the recession windows used in the regime analysis, Table~\ref{tab:regime-agg} reports mean annualized Sharpe ratios by regime, and Figure~\ref{fig:rolling_sharpe} shows the 126-day rolling Sharpe ratio during crises for each strategy.

\begin{table}[!ht]
\caption{NBER recession periods used in the regime analysis}
\centering
\begin{tabular}{@{}ll@{}}
\toprule
Period & Dates \\
\midrule
Dot Com Bubble & 2001-03-01 -- 2001-11-30 \\
Global Financial Crisis & 2007-12-01 -- 2009-06-30 \\
COVID-19 Recession & 2020-02-01 -- 2020-04-30 \\
\bottomrule
\end{tabular}
\label{tab:crisis-periods-combined}
\end{table}

\begin{table}[!ht]
\centering
\caption{Mean annualized Sharpe ratio by regime. Crisis periods are defined using the NBER recession windows in Table~\ref{tab:crisis-periods-combined}, while all remaining observations are classified as calm periods.}
\label{tab:regime-agg}
\resizebox{0.9\textwidth}{!}{%
\begin{tabular}{lrrrrrrr}
\toprule
Regime & HMVA & MVO & HMVA-mv & GMV & HRP & EW & SPY-100 \\
\midrule
Crisis & $-0.055$ & $-0.599$ & $-0.091$ & $-0.213$ & $-0.234$ & $-0.341$ & $-0.184$ \\
Calm   & $0.800$  & $0.280$  & $0.819$  & $0.793$  & $0.790$ & $0.623$ & $0.564$ \\
\bottomrule
\end{tabular}%
}
\end{table}

\begin{figure}[H]
\centering
\includegraphics[width=0.9\linewidth]{results/crisis_analysis/crisis_rolling_sharpe.png}
\caption{126-day rolling Sharpe ratio during each crisis period. Strategy coloring is the same as in the full-sample equity figure.}
\label{fig:rolling_sharpe}
\end{figure}


In the crisis subsample, HMVA records the highest Sharpe ratio, at $-0.055$. HMVA-mv follows with a crisis-period Sharpe ratio of $-0.091$. The remaining crisis-period Sharpe ratios are $-0.184$ for SPY-100, $-0.213$ for GMV, $-0.234$ for HRP, $-0.341$ for EW, and $-0.599$ for MVO.

In the calm subsample, HMVA-mv records the highest Sharpe ratio, at $0.819$. HMVA follows at $0.800$, while GMV and HRP record calm-period Sharpe ratios of $0.793$ and $0.790$, respectively. The passive benchmarks record calm-period Sharpe ratios of $0.623$ for EW and $0.564$ for SPY-100. MVO records the lowest calm-period Sharpe ratio, at $0.280$.


\subsection{Transaction-cost robustness}
\label{sec:cost_robustness}

Table~\ref{tab:cost_robustness_summary} reports after-cost Sharpe ratios for all selected strategies under proportional transaction costs ranging from zero to thirty basis points.

\begin{table}[!ht]
\centering
\caption{The columns $\kappa=0,2,5,10,20,30$ report after-cost Sharpe ratios under proportional  transaction costs measured in basis points.}
\label{tab:cost_robustness_summary}
\begin{tabular}{lrrrrrr}
\toprule
Strategy & $\kappa=0$ & $\kappa=2$ & $\kappa=5$ & $\kappa=10$ & $\kappa=20$ & $\kappa=30$ \\
\midrule
HMVA    & 0.616 & 0.604 & 0.586 & 0.557 & 0.499 & 0.441 \\
MVO     & 0.151 & 0.139 & 0.121 & 0.091 & 0.033 & $-0.025$ \\
\midrule
HMVA-mv & 0.581 & 0.569 & 0.551 & 0.522 & 0.464 & 0.406 \\
GMV     & 0.532 & 0.516 & 0.491 & 0.451 & 0.370 & 0.290 \\
HRP     & 0.498 & 0.492 & 0.482 & 0.466 & 0.433 & 0.400 \\
\midrule
EW      & 0.345 & 0.343 & 0.341 & 0.338 & 0.331 & 0.325 \\
SPY-100 & 0.328 & 0.327 & 0.327 & 0.326 & 0.324 & 0.322 \\
\bottomrule
\end{tabular}
\end{table}

Across the transaction-cost grid, HMVA records the highest after-cost Sharpe ratio in each column. Its Sharpe ratio is $0.616$ at zero transaction costs and declines to $0.441$ at thirty basis points. HMVA-mv starts at $0.581$ for no transaction cost and decreases to $0.406$ when transactions reach 30 basis points. Among the remaining active strategies, GMV declines from $0.532$ to $0.290$, HRP declines from $0.498$ to $0.400$, and MVO declines from $0.151$ to $-0.025$, while the passive benchmarks record smaller numerical changes across the same grid from $0.345$ to $0.325$ for EW and $0.328$ to $0.322$ for SPY-100.

\section{Discussion}
\label{sec:discussion}

The empirical results suggest that HMVA improves point-estimate out-of-sample performance mainly by changing how noisy estimates enter the allocation problem, rather than by eliminating estimation error itself. Across the full sample, HMVA achieves the highest point-estimate Sharpe ratio among the considered strategies, while also recording a materially lower maximum drawdown than most of the benchmarks. This indicates that the improvement is not simply the result of taking more realized risk, but is accompanied by better downside protection in the examined sample.

The comparison with the return-based active benchmark provides useful evidence in favor of HMVA on a point-estimate basis. MVO uses the same estimated inputs as HMVA, but its realized Sharpe ratio is substantially lower and has a deeper maximum drawdown. The Sharpe-ratio difference between HMVA and MVO is positive, although it is not statistically significant after Holm correction within the HMVA hypothesis family. This pattern is consistent with the literature arguing that direct mean-variance optimization can amplify estimation error into unstable and concentrated portfolios. In contrast, HMVA introduces return information only through cluster-level Sharpe proxies and avoids direct covariance-matrix inversion. The positive Sharpe-ratio difference against MVO therefore supports the interpretation that HMVA improves on classical mean-variance optimization in this difficult estimation regime, but the evidence should be interpreted as statistically inconclusive.

HMVA also exceeds both EW and SPY-100 in annualized Sharpe ratio and maximum drawdown performance. This suggests that the proposed hierarchical allocation clears the passive benchmark hurdle on a point-estimate basis in this specification. However, the statistical evidence against the passive benchmarks should be interpreted cautiously. The comparisons with EW and SPY-100 both have Holm-adjusted $p$-values of $0.154$, and the corresponding confidence intervals include zero. Thus, the passive-benchmark evidence is favorable in point estimates, but statistically inconclusive after family-wise correction.

The risk-only version, HMVA-mv, also ranks ahead of the selected risk-only comparators on a point-estimate basis. Its Sharpe ratio exceeds those of GMV and HRP, and its maximum drawdown is below those benchmarks. This suggests that the HMVA tree and recursive allocation rule may add value even without return information. However, the Sharpe-ratio differences relative to GMV and HRP are small and statistically insignificant after Holm correction. The same interpretation applies to its comparison with the passive benchmarks. The point estimates are positive, but the confidence intervals include zero and do not reject equality of Sharpe ratios.

The regime analysis provides further insight into where the performance advantage arises. HMVA weakens during crisis episodes, as expected for a long-only equity strategy, but it weakens less than the competing approaches in the recession windows considered. This suggests that the hierarchical construction is not merely increasing exposure to rewarded equity risk, but is also producing a portfolio that is relatively more robust when correlations rise and diversification becomes more difficult. The risk-only version shows a similar pattern, which indicates that part of the crisis-period resilience comes from the tree structure and recursive risk allocation itself, rather than only from the return signal. During calm periods, HMVA-mv records the strongest performance, while HMVA, GMV, and HRP remain close behind. The regime evidence therefore suggests that HMVA's full-sample advantage is not caused by superiority in every market state, but instead it appears to come from a more balanced profile.

The transaction-cost robustness exercise shows that HMVA's advantage compresses as trading costs increase, which is expected because HMVA and HMVA-mv rebalance more actively than the passive benchmarks. This is an important qualification because part of HMVA's point-estimate advantage is generated by an allocation rule that responds more actively to changing estimated signals and risk structures. The fact that HMVA remains the highest-Sharpe strategy across the tested cost levels suggests that the full-sample advantage is not purely an artifact of the frictionless baseline. At the same time, HRP becomes increasingly competitive at higher transaction costs because its lower turnover makes its performance less sensitive to trading frictions. This highlights a practical trade-off in the proposed method. HMVA uses more information and adapts more aggressively, but this comes at the cost of higher implementation sensitivity. The transaction-cost results should therefore be interpreted as evidence that HMVA is robust to moderate linear costs in this backtest, not as evidence that transaction costs are unimportant.

Taken together, the results support a cautious point-estimate interpretation. HMVA shows positive performance differences relative to classical mean-variance optimization and the passive benchmarks in this specification, but the pairwise Sharpe-ratio tests do not reject equality after Holm correction. The comparisons involving HMVA-mv also remain positive but statistically inconclusive. Therefore, the evidence is best interpreted as showing that HMVA achieves the strongest point-estimate performance among the strategies considered, rather than as conclusive evidence that it dominates every comparator. However, the statistically inconclusive comparisons should not be read as evidence of no effect, since Sharpe-ratio inference is noisy and the effective number of independent observations is limited, thereby limiting the power of the inference tests used here.

\section{Conclusion}
\label{sec:conc}

This thesis set out to construct a mean-variance-inspired portfolio that is less exposed to the estimation error which classical mean-variance optimization amplifies out-of-sample. To that end, it introduced HMVA, a construction that never inverts an estimated covariance matrix and never loads a return vector directly onto individual asset weights. HMVA organizes the asset universe into a top-down binary tree whose splits minimize the correlation between equal-weighted child portfolios. It then allocates capital recursively in proportion to a cluster Sharpe proxy formed from a return estimate, falling back to inverse-volatility bisection wherever the return signal is non-positive to preserve the long-only constraint. A latent weight filter blends each rebalance's target weights with the previously held weights through a data-adaptive gain, reducing trading induced by estimation noise.

Using a fully point-in-time backtest of the top-100 S\&P 500 constituents from July 2000 to January 2024 with a 126-day rolling estimation window, HMVA achieves a full-sample Sharpe ratio of $0.616$, compared with $0.151$ for MVO, $0.345$ for EW, and $0.328$ for SPY-100. Its maximum drawdown is also materially shallower than that of the passive benchmarks. The return-based version improves on the minimum-variance variant, suggesting that the cluster-level Sharpe bisection adds useful information in this sample. The risk-only version, HMVA-mv, also ranks ahead of GMV and HRP on point estimates, but these differences are not statistically significant after correction. 

The research question was whether mean and variance estimates can be introduced into an allocation framework to outperform both naive and sophisticated benchmarks in out-of-sample risk-adjusted returns, given identical inputs in a difficult estimation regime. The empirical answer is cautiously positive on point estimates, but statistically conservative. In the setting studied here, HMVA achieves the highest full-sample point-estimate Sharpe ratio among the strategies considered and records positive Sharpe-ratio differences against all comparators in the HMVA hypothesis family. It also remains the highest-Sharpe strategy across the tested transaction-cost levels. However, the pairwise Sharpe-ratio tests do not reject equality after Holm correction and should therefore be interpreted cautiously, while the comparisons involving HMVA-mv are positive but statistically inconclusive. The regime analysis further suggests that HMVA's full-sample performance is supported by relatively strong crisis-period performance together with solid calm-period performance, rather than by dominance in every regime. The results should therefore be interpreted as evidence that HMVA performs best on point estimates in this specification, rather than as conclusive evidence that it dominates every comparator.

Several limitations remain. The evaluation covers a single market and asset class, namely large-capitalization US equities, so the signal and correlation structures HMVA exploits may behave differently in international, small-capitalization, or multi-asset universes. Although the contiguous-cut split heuristic tracks the exhaustive optimum closely in objective value, it is validated on the split objective rather than on realized portfolio performance. As with any newly proposed strategy, some degree of in-sample design bias may also remain despite the point-in-time backtest, since the method was developed and evaluated on the same broad empirical setting. Finally, HMVA contains several modeling choices, including the approximate split threshold, the exponential half-life, the cluster moment proxies, and the other uncertainty proxies used in the allocation, whose sensitivity and optimality have not been fully explored.

These limitations indicate natural directions for future work. Extending the evaluation to the full S\&P 500, to international equity markets, and to other asset classes would test whether the construction retains its observed advantage in settings where estimation error is more severe. Longer samples and broader datasets would also help determine whether the observed Sharpe differences can be established with stronger statistical evidence. Future work could also evaluate HMVA under more realistic institutional constraints, including liquidity, capacity, sector exposure limits, and nonlinear transaction costs. Finally, replacing the entropy-and-velocity heuristic proxies with more explicit uncertainty measures for the adaptive latent weight filter, and testing alternative return forecasts such as factor-based or machine-learning signals, could improve stability and performance around structural breaks while preserving the central idea of hierarchical, estimation-error-resistant mean-variance allocation.


\bibliographystyle{chicago}
\bibliography{references}

\newpage
\section{Appendix}

\subsection{Implementation}
\label{app:pseudo-code}
The three conceptual stages of HMVA described in Section~\ref{sec:hmva} translate almost directly into computation, with the main exception being the minimization of the split objective $f(L,R)$. This objective cannot be evaluated exactly at scale, because a cluster of size $n$ admits $2^{n-1}-1$ distinct bipartitions. This is trivial near the leaves but intractable near the root, where $n$ approaches the full universe size. The algorithm therefore resolves the split with one of two routines chosen by a size threshold $n_{\text{bf}}$, set to $10$ throughout in the thesis implementation. The algorithm below presents the full pseudo-code procedure at one rebalance date.

\vspace{0.5 cm}

\begin{algorithmic}[1]
\toprule
\REQUIRE return window $R_{t-T_h:t-1}$, previous weights $w_{old}$, prior signal $\hat{\mu}_{old}$, threshold $n_{\text{bf}}$
\midrule
\STATE $\hat{\Sigma}_t \gets \hat{\Sigma}(R_{t-T_h:t-1})$ \hfill \COMMENT{(\ref{eq:nls-cov})}
\STATE $\hat{\mu}_{t,new}\gets \hat{\mu}(R_{t-T_h:t-1},\hat{\Sigma}_t)$ \hfill \COMMENT{(\ref{eq:BL_simplified})}
\STATE $\textbf{root} \gets$ cluster of all $N$ assets, $\textbf{queue} \gets \{\textbf{root}\}$
\WHILE{$\textbf{queue}$ is non-empty}
    \STATE pop cluster $C$ from $\textbf{queue}$
    \IF{$|C| \le 1$}
        \STATE mark $C$ as a leaf, \textbf{continue}
    \ENDIF
    \IF{$|C| \le n_{\text{bf}}$}
        \STATE $(L,R) \gets \arg\min f(\cdot,\cdot)$ over all subsets \hfill \COMMENT{exhaustive}
    \ELSE
        \STATE sort $C$ by within-cluster row-sum
        \STATE $(L,R) \gets$ best contiguous cut via prefix sums \hfill \COMMENT{approximate}
    \ENDIF
    \STATE attach $L,R$ as children of $C$, push $L,R$ onto $\textbf{queue}$
\ENDWHILE
\STATE \textbf{Allocate}$(\textbf{root},\text{budget}=1)$ recursively
\STATE \quad compute $\hat{S}(L),\hat{S}(R)$ at the node \hfill \COMMENT{(\ref{eq:sharpe-proxy})}
\STATE \quad \textbf{if} $\hat{S}(L)+\hat{S}(R)>0$ \textbf{ then } $\alpha_L \gets \hat{S}(L)/(\hat{S}(L)+\hat{S}(R))$
\STATE \quad \textbf{else } $\alpha_L \gets \widehat{IV}(L)/(\widehat{IV}(L)+\widehat{IV}(R))$ \hfill \COMMENT{(\ref{eq:inv_var})}
\STATE \quad send $\text{budget } \alpha_L$ to $L$ and $\text{budget }(1-\alpha_L)$ to $R$
\STATE \quad recurse until each leaf holds a single weight, collect into $w_{new}$
\STATE $K \gets Y/(X+Y)$
\STATE $\tilde{w}_t \gets K w_{new} + (1-K)w_{old}$
\STATE \textbf{return} normalized weights $\tilde{w}_t / \sum_i \tilde{w}_{t,i}$
\end{algorithmic}
\bottomrule

\vspace{0.5 cm}

For any cluster with $n_C\le n_{\text{bf}}$, the bipartition is found by exhaustive search. Exploiting the $L\!\leftrightarrow\!R$ symmetry of $f$, it suffices to enumerate the subsets of size $1,\dots,\lfloor n/2\rfloor$ and assign the complement to the other branch, evaluating the equal-weighted cross-cluster correlation $f(L,R)$ for each and keeping the minimizer. Each candidate split is scored using the cross-block covariance sum and the two within-block variance sums, which together determine the equal-weighted correlation in \eqref{eq:split-obj}. The exhaustive search therefore returns the global optimum of the split objective for these small clusters.

For $n_C>n_{\text{bf}}$, exhaustive enumeration is replaced by an approximate heuristic that restricts the search to contiguous cuts of a single ordering. The assets of the cluster are first sorted by their within-cluster covariance row-sum $\sum_{j\in C}\hat{\Sigma}_{ij}$, which places strongly co-moving assets adjacent in the sequence. A two-dimensional prefix-sum array of the sorted covariance block is then formed once in $O(n^2)$, after which the cross-block and within-block sums entering the correlation objective can be read off in constant time for each of the $n-1$ possible cuts.

To justify the use of the approximation, a simulation is conducted on $100$ independently generated random clusters at each size $n\in\{12,14,16,18,20\}$, giving $500$ clusters in total. These are the largest sizes at which exhaustive search remains feasible. For each cluster, the heuristic split is compared against the exhaustive optimum of the objective $f(L,R)$. Figure~\ref{fig:split-quality} plots the heuristic objective against the exhaustive objective. Points on the line are clusters where the heuristic attains the optimum, while points above it correspond to heuristic splits with weaker objective values than the optimum.

\begin{figure}[H]
\centering
\includegraphics[width=0.66\linewidth]{results/split_comparison/score_scatter.png}
\caption{Heuristic versus exhaustive split objective across $500$ random clusters ($100$ each at sizes $n\in\{12,14,16,18,20\}$). Each point is one cluster, plotting the contiguous-cut heuristic objective against the exhaustive optimum. The dashed line marks $y=x$, where the two agree exactly.}
\label{fig:split-quality}
\end{figure}

\begin{table}[H]
\centering
\begin{tabular}{lrrrr}
\toprule
$n$ & Median ratio & 95th pct.\ ratio & Exact match \\
\midrule
12 & 1.000 & 1.000 & 71.0\% \\
14 & 1.000 & 1.000 & 69.0\% \\
16 & 1.000 & 1.000 & 69.0\% \\
18 & 1.000 & 1.000 & 69.0\% \\
20 & 1.000 & 1.000 & 56.0\% \\
\bottomrule
\end{tabular}
\caption{Heuristic split quality relative to exhaustive search. Values equal to one indicate that the heuristic attains the exhaustive objective value.}
\label{tab:split-quality}
\end{table}
The code and implementation files used in this thesis, excluding the CRSP data, are available at
\text{\url{https://github.com/Developer-KI/BSc_Thesis}}.

\subsection{Benchmarks}
\label{app:benchmarks}
This appendix specifies the allocation used by each competing strategy. All return-based and variance-based optimizers receive the same estimated inputs $\hat{\Sigma}$ and $\hat{\mu}$ and operate on the same point-in-time asset universe, and are constrained to be long-only and fully invested in the universe at all times.

\noindent \textbf{MVO}: The classical Markowitz utility maximizer trades expected return against variance under a risk-aversion coefficient $\gamma$, set to $\gamma = 2.5$ throughout:
\begin{equation}
  w^{\text{MVO}} = \arg\max_{w}\;
    w^{\top}\hat{\mu} - \tfrac{\gamma}{2}\, w^{\top}\hat{\Sigma}\, w
  \qquad \text{s.t.}\quad
  \mathbf{1}^{\top} w = 1,\; 0 \le w_i \le 1
  \label{eq:app-mvo}
\end{equation}
The problem is solved by sequential least-squares quadratic programming, with the equal-weight portfolio used as the starting point and as a fallback on solver failure. Before optimization the covariance is symmetrized and its eigenvalues are clipped to a small positive floor to guarantee positive definiteness.

\noindent \textbf{GMV}: The minimum-variance portfolio discards the return estimate entirely and minimizes portfolio variance subject to the same constraints:
\begin{equation}
  w^{\text{GMV}} = \arg\min_{w}\;
    w^{\top}\hat{\Sigma}\, w
  \qquad \text{s.t.}\quad
  \mathbf{1}^{\top} w = 1,\; 0 \le w_i \le 1 
  \label{eq:app-gmv}
\end{equation}
again solved by sequential least-squares quadratic programming from an equal-weight start. Before optimization the covariance is symmetrized and its eigenvalues are clipped to a small positive floor to guarantee positive definiteness.

\noindent \textbf{HRP}: This benchmark converts the estimated correlation matrix to the distance metric $d_{ij} = \sqrt{\tfrac{1}{2}(1 - \rho_{ij})}$, on which single-linkage agglomerative clustering produces an asset ordering via quasi-diagonalization. Capital is then assigned by recursive bisection of this ordered list. At each split, the two subclusters $L$ and $R$ receive weights in proportion to their inverse-volatility cluster measures:
\begin{equation}
\alpha_L = \frac{\widehat{IV}(L)}{\widehat{IV}(L) + \widehat{IV}(R)},
\quad
\alpha_R= 1 - \alpha_L
\end{equation}

\noindent \textbf{SPY-100}: The passive value-weighted benchmark allocates in proportion to point-in-time market capitalization:
\begin{equation}
  w_i^{\text{SPY-100}} = \frac{\text{cap}_i}{\sum_{\text{all n}} \text{cap}_n}
  \label{eq:app-spy}
\end{equation}
using the most recent available daily market-cap figure for each constituent of the top-100 universe at the rebalance date, using neither $\hat{\Sigma}$ nor $\hat{\mu}$. It falls back to equal weighting on any rebalance where cap data is unavailable.

\noindent \textbf{EW}: The naive $1/N$ benchmark assigns $w^\mathrm{EW}_i = 1/N$ to every asset in the universe at each rebalance, using neither $\hat{\Sigma}$ nor $\hat{\mu}$. It serves as the minimum hurdle that any active strategy must clear.

\end{document}
